\metatitle{Series-parallel posets and stable lifts}
\metadescription{A recursive colored-matroid encoding of linear extensions
of series-parallel posets, with descents counted by one color and real-stable marginals.}

\section[seriesParallel]{Series-parallel posets and stable lifts}

Series-parallel posets have real-rooted $P$-Eulerian polynomials
\cite{Wagner1992en,Wagner1992}. We strengthen this scalar statement by
constructing a colored matroid whose words are in bijection with the linear
extensions, with descents counted by one color on marked positions.
Its color marginal is real stable. The construction gives both an exchange
structure on the extensions and a multivariate explanation of real-rootedness.

\subsection[seriesParallelPosets]{Posets and descents}

A \defin{series-parallel poset} is obtained from singletons by disjoint
unions and ordinal sums. In $P\sqcup Q$ there are no comparisons between
the two parts; in $P\oplus Q$ every element of $P$ is below every element
of $Q$. A rooted tree in ancestor order is its root followed by the disjoint
union of its child subtrees, so rooted forests belong to this class.

We use natural labelings: $i<_P j$ implies $i<j$ as integers.
A \defin{linear extension} is a permutation $\pi$ compatible with $P$.
Write $\mathcal L(P)$ for the extensions, $h(P)$ for the number of elements
in a longest chain, and
\[
 \begin{gathered}
 \operatorname{des}(\pi)
 =|\{i:\pi_i>\pi_{i+1}\}|,\\
 W_P(t)=\sum_{\pi\in\mathcal L(P)}t^{\operatorname{des}(\pi)}.
 \end{gathered}
\]

\subsection[seriesParallelLift]{The colored-matroid lift}

For a word $u$ on a position set $E$, put
$B(u,c)=\{e:u(e)=c\}$. A nonzero real polynomial is \defin{real stable}
if it does not vanish when all its variables have positive imaginary part.

\begin{theorem}[Series-parallel posets]
Let $P$ be a naturally labeled series-parallel poset with $n$ elements.
There are allowed-color lists $L_e$ on a finite position set $E$, fixed
color quotas, and a set $S\subseteq E$ of marked positions such that the
family $\mathcal F$ of all words satisfying the lists and quotas admits
a bijection $\pi\mapsto u_\pi$ from $\mathcal L(P)$ with
\[
 \begin{gathered}
 |B(u_\pi,1)|=n-h(P),\\
 |B(u_\pi,1)\cap S|=\operatorname{des}(\pi).
 \end{gathered}
\]
The family $\mathcal F$ is a colored matroid. Its marginal
\[
 f(\mathbf x)=\sum_{u\in\mathcal F}\prod_{e\in B(u,1)}x_e
\]
is homogeneous of degree $n-h(P)$ and real stable. Setting $x_e=t$ for
$e\in S$ and $x_e=s$ otherwise gives
\[
 \sum_{\pi\in\mathcal L(P)}
 t^{\operatorname{des}(\pi)}s^{n-h(P)-\operatorname{des}(\pi)}.
\]
\end{theorem}

\begin{proof}
The recursive construction below gives the lists, quotas, bijection and
marked statistic. The \href{constructions.htm#allowedColors}{allowed-color
proposition} proves symmetric cyclic exchange. The stability proposition
below applies to these lists.
\end{proof}

Thus $W_P(t)$ is obtained by setting the marked variables equal to $t$
and the others to $1$. Real specialization and identification of variables
preserve stability, so $W_P$ is real-rooted. Its constant term is $1$ and
its coefficients are nonnegative, hence every zero is negative.

\begin{example}[A rooted tree]
Let $P$ have relations generated by $1<2<3$ and $1<4<5$.
The construction gives six positions, colors $1,2,3,4$, and quotas
$(2,2,1,1)$. The lists and marks are
\[
\begin{array}{c|c|c}
 e&L_e&e\in S\\ \hline
 1&\{1,2\}&-\\
 2&\{1,2\}&-\\
 3&\{2,3\}&-\\
 4&\{1,3\}&\checkmark\\
 5&\{2,4\}&-\\
 6&\{1,4\}&\checkmark
\end{array}
\]
All permitted words appear in the following bijection.
\[
\begin{array}{c|c|c}
 \pi&u_\pi&\operatorname{des}(\pi)\\ \hline
 \mathtt{12345}&\mathtt{112324}&0\\
 \mathtt{12435}&\mathtt{123124}&1\\
 \mathtt{12453}&\mathtt{122341}&1\\
 \mathtt{14235}&\mathtt{213124}&1\\
 \mathtt{14253}&\mathtt{223141}&2\\
 \mathtt{14523}&\mathtt{212341}&1
\end{array}
\]
Between $\mathtt{112324}$ and $\mathtt{223141}$, positions $1,3,4$
form the color cycle $1\to2\to3\to1$. Exchanging those positions gives
$\mathtt{213124}$ and $\mathtt{122341}$, both in the table.
The marginal is
\[
 \begin{aligned}
 f(\mathbf x)&=x_1x_2+x_1x_4+x_1x_6\\
 &\quad+x_2x_4+x_2x_6+x_4x_6.
 \end{aligned}
\]
Setting $x_1=x_2=s$ and $x_4=x_6=t$ gives $s^2+4st+t^2$.
Hence $W_P(t)=1+4t+t^2$.
\end{example}

\subsection[listStability]{Stability from allowed colors}

The stability argument works for every allowed-color family, not just the
ones constructed from posets. It also permits positive weights on the
assignments in the other colors.

\begin{proposition}[Weighted color marginals]
Let $\mathcal F$ be a nonempty family of all words with lists $L_e$ and
quotas $r_1,\dotsc,r_k$. Give each allowed assignment $(e,c)$ with $c\ne1$
a positive weight $a_{e,c}$. Then
\[
 f(\mathbf x;\mathbf a)
 =\sum_{u\in\mathcal F}
 \left(\prod_{e:u(e)\ne1}a_{e,u(e)}\right)
 \prod_{e:u(e)=1}x_e
\]
is homogeneous of degree $r_1$ and real stable.
\end{proposition}

\begin{proof}
Introduce a variable $z_c$ for every color $c\ne1$. Define one linear
form for each position by
\[
 \ell_e=\mathbf1_{1\in L_e}x_e+
 \sum_{c\in L_e\setminus\{1\}}a_{e,c}z_c.
\]
Expanding their product chooses one allowed color at each position.
Consequently
\[
 f(\mathbf x;\mathbf a)
 =\left[\prod_{c=2}^k z_c^{r_c}\right]\prod_{e\in E}\ell_e.
\]
The extracted quotas force the remaining quota $r_1$.
Each linear form is stable, since its coefficients are nonnegative and
not all zero. Their product is stable. Coefficient extraction is repeated
differentiation followed by specialization at $z_c=0$, divided by the
appropriate factorials. These operations preserve stability or give zero;
nonemptiness and positive weights exclude zero here. See
\cite[Section~2.1 and Lemma~2.9]{BrandenHuh2020} for the stability properties.
\end{proof}

Every word contributes to the marginal, so the coefficient of a target-set
monomial counts all its completions in the other colors. The proposition
applies with any color distinguished. Since homogeneous real-stable
polynomials with nonnegative coefficients are Lorentzian
\cite[Proposition~2.2]{BrandenHuh2020}, this verifies a stronger form of the
\href{rank.htm#lorentzianConjecture}{completion-polynomial conjecture}
for allowed-color families.

\subsection[seriesParallelRecursion]{The recursive construction}

We maintain a bijection to all words satisfying lists and quotas, with
color $1$ counting descents on a marked set $S$. Put
$D=\{e:1\in L_e\}$ and let $r$ be the color-$1$ quota.
For a poset $P$, the additional conditions are
\[
 \begin{gathered}
 r=|P|-h(P),\qquad |D\setminus S|=r,\\
 |S|\leq|P|,\qquad S\subseteq D.
 \end{gathered}
\]
The positions are auxiliary: they need not be elements of the poset.
We keep the declared lists, including permissions for colors of quota zero.

\subsubsection[quotaReversal]{Reversing the quotas}

The following operation exchanges the numbers of marked and unmarked
declared color-$1$ positions while preserving the statistic.

\begin{lemma}[Quota reversal]
Suppose the color-$1$ quota is $r$, with $r$ unmarked and $M$ marked
positions allowing color $1$. There is a bijection to another allowed-color
family with quota $M$, $M$ unmarked and $r$ marked positions allowing color
$1$, preserving the marked color-$1$ count.
\end{lemma}

\begin{proof}
Replace each position $e$ by positions $(e,c)$ for $c\in L_e$.
Introduce a fresh color $p_e$ of quota one and give $(e,c)$ the list
$\{c,p_e\}$. For an old word $u$, put
\[
 v(e,c)=\begin{cases}p_e&u(e)=c,\\c&u(e)\ne c.\end{cases}
\]
If $c$ was allowed at $d_c$ positions and had quota $r_c$, give it new
quota $d_c-r_c$. The unique occurrence of $p_e$ recovers $u(e)$;
all other cells in that row retain their old color. The complemented
quotas recover the old quotas, so the image is the entire new family.

Mark $(e,1)$ precisely when $e$ was unmarked and allowed color $1$.
If the old marked count was $j$, its unmarked count was $r-j$.
The new marked count is $r-(r-j)=j$. The new quota and the numbers of
marked and unmarked declared positions are as stated.
\end{proof}

\subsubsection[shuffleBijection]{The shuffle bijection}

Let $\sigma,\tau$ be permutations on disjoint alphabets of sizes $n,m$,
with every letter of $\sigma$ smaller than every letter of $\tau$.
Put $k=\operatorname{des}(\sigma)$ and $\ell=\operatorname{des}(\tau)$.
For disjoint ordered sets $U,V$ with
\[
 |U|=n-k+\ell,\qquad |V|=m-\ell+k,
\]
I.~P.~Goulden's
\href{https://uwaterloo.ca/math/sites/default/files/uploads/documents/gtams1985.pdf}{
bijective proof of Stanley's shuffling theorem}
(Theorems~3.4--3.5) gives a bijection between shuffles $\pi$ of
$\sigma,\tau$ and $n$-subsets $R\subseteq U\sqcup V$ satisfying
\[
 \operatorname{des}(\pi)=\ell+|R\cap V|.
\]
In Goulden's formulation, a shuffle with $d$ descents gives subsets
$\alpha\subseteq U$, $\beta\subseteq V$ of sizes $d-k$, $d-\ell$.
Take $R=(U\setminus\alpha)\sqcup\beta$. Conversely,
$\alpha=U\setminus R$ and $\beta=V\cap R$ recover the shuffle through
his inverse. We fix this bijection and the orders on $U,V$ throughout.

\subsubsection[seriesParallelBase]{Singletons and ordinal sums}

A singleton has one extension. Encode it by the empty word, with color-$1$
quota zero. For $P\oplus Q$, take disjoint position sets and disjoint
palettes except for color $1$. Concatenate the parent words, add the
color-$1$ quotas, and retain their marks and other quotas.

The other colors force the color-$1$ count separately in each part,
so splitting the word gives the inverse on the whole allowed-color family.
Label $P$ before $Q$. Extensions concatenate without a boundary descent,
and heights add. The required statistic and all three numerical conditions
are therefore preserved.

\subsubsection[seriesParallelUnion]{Disjoint unions}

Write $n=|P|$, $m=|Q|$, $r_P=n-h(P)$ and $r_Q=m-h(Q)$.
Let $M_P,M_Q$ be their numbers of marked declared color-$1$ positions.
Interchange the parents if necessary so that $h(Q)\geq h(P)$, and label
$P$ before $Q$. Keep their other colors disjoint.

\begin{proposition}[Disjoint-union construction]
From encodings satisfying the stated conditions for $P$ and $Q$, we obtain
one for $P\sqcup Q$ with color-$1$ quota $n+r_Q$.
\end{proposition}

\begin{proof}
We first prepare the two parent words, then encode their shuffle by a
choice of colors, and finally apply quota reversal to count its descents.

Apply quota reversal to $Q$. Add $n-r_P$ unmarked positions forced to
color $1$ in $P$, and $m-M_Q$ such positions in the reversed $Q$.
Every padded word now has $n$ or $m$ occurrences of color $1$,
respectively. The unmarked declared counts are $n,m$, and the marked
declared counts are $M_P,r_Q$.

Allow a fresh color $2$ wherever color $1$ is allowed, renaming other
colors first if needed. Give colors $1,2$ global quotas $n,m$ and keep
the other quotas. Merging $1,2$ recovers the two padded parent words.
Thus the whole new family consists of a parent pair and an $n$-subset
$R$ of its occupied target positions, colored $1$; the rest are colored $2$.

For parents with $k,\ell$ descents, let $U$ consist of the unmarked
occupied targets in $P$, followed by the marked occupied targets in $Q$.
Let $V$ consist of the marked occupied targets in $P$, followed by the
unmarked occupied targets in $Q$. Fix an order within each group.
Their sizes are $n-k+\ell,m-\ell+k$, so the shuffle bijection applies.
Its descent count is $\ell$ plus the number of selected positions in $V$:
the marked $P$ entries and unmarked $Q$ entries colored $1$.

Leave the $P$ part unchanged. On the padded $Q$ part apply the row
replacement from quota reversal, but interchange retained colors $1,2$.
Explicitly, for its current lists $K_e$, let $\rho$ swap $1,2$ and fix
other colors. Replace $e$ by cells $(e,c)$, $c\in K_e$, with lists
$\{\rho(c),p_e\}$ and fresh colors $p_e$ of quota one. An entry $u(e)$
puts $p_e$ at $(e,u(e))$ and $\rho(c)$ in the other cells.
Complement each old hidden-color quota within the $Q$ part.
There are $d=m+r_Q$ declared target positions in that part. Put
\[
 r_1=n+r_Q,\qquad r_2=d.
\]
Keep the $P$ marks, and mark $(e,2)$ in the new $Q$ part exactly when
$e$ was an unmarked declared target position.

To verify the full image and inverse, let $a_Q,b_Q$ be the old numbers
of colors $1,2$ in the $Q$ part. Then $a_Q+b_Q=m$, and their new counts
are $d-b_Q,d-a_Q$. If $a_P$ is the old color-$1$ count in $P$, the
old condition $a_P+a_Q=n$ gives the displayed new quotas.
Conversely, the unique $p_e$ in each row recovers every old $Q$ entry.
The complemented hidden quotas force $a_Q+b_Q=m$; the unchanged hidden
quotas in $P$ force its combined $1,2$ count to be $n$.
The final color-$1$ quota gives
\[
 a_P+d-b_Q=n+r_Q=n+d-m,
\]
hence $a_P+a_Q=n$. We have recovered exactly a word from the selection
step. Merge colors $1,2$, remove the padding, and invert the initial
reversal to recover the parent extensions. They recover $U,V$, and the
old color-$1$ positions recover $R$; Goulden's inverse recovers the shuffle.

There were $m$ unmarked declared target positions in $Q$. Of these,
$m-\ell$ carried colors $1,2$ and $\ell$ carried other colors.
The final marked color-$1$ positions in this part correspond to all old
unmarked positions except those colored $2$. Their number is therefore
$\ell$ plus the number of old unmarked $Q$ entries colored $1$.
Together with the $P$ marks, this is the shuffle's descent number.
The final unmarked declared count is $n+r_Q$, and the marked declared
count is $M_P+m\leq n+m$. Finally,
\[
 \begin{aligned}
 n+r_Q&=n+m-h(Q)\\
 &=|P\sqcup Q|-h(P\sqcup Q).
 \end{aligned}
\]
\end{proof}

For a prescribed natural labeling, transport the construction along adjacent
swaps of incomparable elements in the natural label order. When swapping
consecutive labels $j,j+1$, leave an extension's numerical word unchanged
if those labels are adjacent in it; otherwise interchange their values.
This is a descent-preserving bijection to extensions in the new labeling,
inverted by the same rule. Any two natural label orders are connected by
such swaps. This completes the construction for every natural labeling.

\subsection[stableLiftComparison]{Comparison with multivariate Eulerian polynomials}

P.~Br\"and\'en and M.~Leander \cite{BrandenLeander2016x} study the
refinement
\[
 A_P(\mathbf x,\mathbf y)
 =\sum_{\pi\in\mathcal L(P)}
 \prod_{i\in\operatorname{DB}(\pi)}x_i
 \prod_{j\in\operatorname{AB}(\pi)}y_j,
\]
where $\operatorname{DB}(\pi)$ and $\operatorname{AB}(\pi)$ are the sets
of descent and ascent bottoms, with boundary values
$\pi_0=\pi_{n+1}=\infty$. Thus $\pi_i$ is a descent bottom when
$\pi_{i-1}>\pi_i$, and an ascent bottom when $\pi_i<\pi_{i+1}$.
Their polynomial is homogeneous of degree $n+1$ and specializes to
$tW_P(t)$ at $\mathbf x=t$, $\mathbf y=1$.

They prove stability under disjoint union (Theorem~2.2 and Corollary~2.3)
and for naturally labeled decreasing forests (Corollary~5.5). Their
ordinal-sum lemma assumes stability for $P$ and for $Q$ with a new minimum
adjoined (Lemma~5.4).

Our variables instead record positions in the colored-word construction.
The resulting lift has degree $n-h(P)$ and is built recursively for every
series-parallel poset. It retains a bijection to complete colorings and
symmetric cyclic exchange, while every positively weighted color marginal
is stable. This strengthens the classical real-rootedness statement by
exhibiting the combinatorial structure behind a stable lift.
